\chapter{Conclusioni}
\label{cha:conclusioni}
Dal confronto tra i due modelli emergono alcune differenze interessanti. L'SVC ha ottenuto risultati più limitati, ma permette di ricondurre la classificazione a un'equazione quadratica, rendendo esplicita la frontiera decisionale e quindi più facile da interpretare. L'MLP ha invece fornito risultati molto buoni: osservando la sua frontiera decisionale, si nota che ricalca molto da vicino la distribuzione degli esiti già visibile nelle mappe di calore della Sezione~\ref{subsec:relazione-parametri-esiti}.
Le difficoltà maggiori si incontrano nella zona di transizione tra sopravvivenza e dissoluzione, dove ammassi con gli stessi parametri orbitali possono avere esiti diversi a seconda della loro realizzazione iniziale. Questo potrebbe indicare che semiasse maggiore ed eccentricità, da soli, non siano sufficienti a descrivere il risultato. Sarebbe quindi interessante, in futuro, aumentare il numero di realizzazioni per ogni configurazione orbitale e includere altre grandezze diagnostiche tra i dati in ingresso, così da fornire ai modelli informazioni più complete.
Un'altra difficoltà emersa durante lo studio riguarda la dissoluzione autonoma degli ammassi, dovuta alla loro stessa evoluzione interna e non direttamente all'azione del potenziale galattico. Proprio per limitarne l'influenza, il tempo finale delle simulazioni è stato fissato a $T_{\mathrm{end}}=700$. Anche questo aspetto potrebbe essere approfondito in futuro, ad esempio aumentando il numero di stelle utilizzato per ciascun ammasso.
Un risultato positivo riguarda invece il modello di Plummer utilizzato per generare le condizioni iniziali. Sebbene la sua descrizione teorica sia formulata nel limite $N\to\infty$, si è dimostrato efficace anche con ammassi di sole 200 stelle. Il riscalamento delle velocità iniziali si è rivelato particolarmente importante per ottenere configurazioni in equilibrio viriale.
Infine, l'intervallo dei semiassi maggiori esplorato costituisce un altro limite del dataset. Sarebbe interessante ampliarlo e verificare se i modelli riescano a prevedere l'esito di configurazioni con parametri orbitali al di fuori degli intervalli utilizzati durante l'addestramento. In questo modo si potrebbe valutare la loro capacità di estrapolare le relazioni apprese e applicarle anche a configurazioni mai simulate in precedenza.